\section{Enfoque de la investigación}

Este proyecto utiliza un enfoque cuantitativo, que se basa en recoger y analizar datos numéricos para describir fenómenos a través de métodos estadísticos y matemáticos \cite{cuantitativa}. Este enfoque prioriza la objetividad y medición precisa para relacionar causales entre las variables observadas.

A diferencia de los enfoques cualitativos o mixtos, este no requiere de estudiar percepciones, significados ni experiencias humanas, por lo que es adecuado para investigaciones donde las variables pueden ser definidas, controladas y medidas \cite{mixto}.

El enfoque cuantitativo se justifica porque el análisis se realiza a partir de datos generados mediante simulación computacional. En este proceso no hay intervención directa con humanos y los resultados obtenidos no están sujetos a interpretaciones. Este enfoque permite comparar escenarios, desempeño de modelos e influencia de parámetros sobre la dinámica emergente del sistema.

